\documentclass[10pt,a4paper]{amsart}
\usepackage{amsmath,amssymb,amsthm,graphicx}
\usepackage{verbatim}
\newenvironment{reference}{\comment}{\endcomment}
%\newenvironment{reference}{}{}
\newenvironment{slogan}{\comment}{\endcomment}
\newenvironment{history}{\comment}{\endcomment}

% For commutative diagrams we use Xy-pic
\usepackage[all]{xy}

% We use 2cell for 2-commutative diagrams.
\xyoption{2cell}
% Theorem environments.
%
\theoremstyle{plain}
\newtheorem{theorem}[subsection]{Theorem}
\newtheorem{proposition}[subsection]{Proposition}
\newtheorem{lemma}[subsection]{Lemma}

\theoremstyle{definition}
\newtheorem{definition}[subsection]{Definition}
\newtheorem{example}[subsection]{Example}
\newtheorem{exercise}[subsection]{Exercise}
\newtheorem{situation}[subsection]{Situation}

\theoremstyle{remark}
\newtheorem{remark}[subsection]{Remark}
\newtheorem{remarks}[subsection]{Remarks}

\numberwithin{equation}{subsection}

% Macros
%
\def\lim{\mathop{\mathrm{lim}}\nolimits}
\def\colim{\mathop{\mathrm{colim}}\nolimits}
\def\Spec{\mathop{\mathrm{Spec}}}
\def\Hom{\mathop{\mathrm{Hom}}\nolimits}
\def\Ext{\mathop{\mathrm{Ext}}\nolimits}
\def\SheafHom{\mathop{\mathcal{H}\!\mathit{om}}\nolimits}
\def\SheafExt{\mathop{\mathcal{E}\!\mathit{xt}}\nolimits}
\def\Sch{\mathit{Sch}}
\def\Mor{\mathop{\mathrm{Mor}}\nolimits}
\def\Ob{\mathop{\mathrm{Ob}}\nolimits}
\def\Sh{\mathop{\mathit{Sh}}\nolimits}
\def\NL{\mathop{N\!L}\nolimits}
\def\CH{\mathop{\mathrm{CH}}\nolimits}
\def\proetale{{pro\text{-}\acute{e}tale}}
\def\etale{{\acute{e}tale}}
\def\QCoh{\mathit{QCoh}}
\def\Ker{\mathop{\mathrm{Ker}}}
\def\Im{\mathop{\mathrm{Im}}}
\def\Coker{\mathop{\mathrm{Coker}}}
\def\Coim{\mathop{\mathrm{Coim}}}

% Boxtimes
%
\DeclareMathSymbol{\boxtimes}{\mathbin}{AMSa}{"02}

%
% Macros for moduli stacks/spaces
%
\def\QCohstack{\mathcal{QC}\!\mathit{oh}}
\def\Cohstack{\mathcal{C}\!\mathit{oh}}
\def\Spacesstack{\mathcal{S}\!\mathit{paces}}
\def\Quotfunctor{\mathrm{Quot}}
\def\Hilbfunctor{\mathrm{Hilb}}
\def\Curvesstack{\mathcal{C}\!\mathit{urves}}
\def\Polarizedstack{\mathcal{P}\!\mathit{olarized}}
\def\Complexesstack{\mathcal{C}\!\mathit{omplexes}}
% \Pic is the operator that assigns to X its picard group, usage \Pic(X)
% \Picardstack_{X/B} denotes the Picard stack of X over B
% \Picardfunctor_{X/B} denotes the Picard functor of X over B
\def\Pic{\mathop{\mathrm{Pic}}\nolimits}
\def\Picardstack{\mathcal{P}\!\mathit{ic}}
\def\Picardfunctor{\mathrm{Pic}}
\def\Deformationcategory{\mathcal{D}\!\mathit{ef}}
\usepackage{cleveref}
\crefname{originalreference}{original reference}{original references}
\Crefname{originalreference}{Original reference}{Original references}
\makeatletter
\newlabel{situation-versal}{{0DQU}{source}{}{sourcepointer-1}{}}
\newlabel{situation-versal@cref}{{[originalreference][0][]0DQU}{[1][0][]source}{}{}{}}
\makeatother
\crefname{originalremark}{original remark}{original remarks}
\Crefname{originalremark}{Original remark}{Original remarks}
\makeatletter
\newlabel{remark-upgrade}{{0DR1}{source}{}{sourcepointer-2}{}}
\newlabel{remark-upgrade@cref}{{[originalremark][0][]0DR1}{[1][0][]source}{}{}{}}
\makeatother
\crefname{originallemma}{original lemma}{original lemmas}
\Crefname{originallemma}{Original lemma}{Original lemmas}
\makeatletter
\newlabel{lemma-Artin-approximation-by-smooth-morphism}{{0DR0}{source}{}{sourcepointer-3}{}}
\newlabel{lemma-Artin-approximation-by-smooth-morphism@cref}{{[originallemma][0][]0DR0}{[1][0][]source}{}{}{}}
\makeatother
\crefname{originallemma}{original lemma}{original lemmas}
\Crefname{originallemma}{Original lemma}{Original lemmas}
\makeatletter
\newlabel{lemma-versal-ring}{{0DQX}{source}{}{sourcepointer-4}{}}
\newlabel{lemma-versal-ring@cref}{{[originallemma][0][]0DQX}{[1][0][]source}{}{}{}}
\makeatother
\crefname{originalreference}{original reference}{original references}
\Crefname{originalreference}{Original reference}{Original references}
\makeatletter
\newlabel{stacks-morphisms-lemma-composition-flat}{{06PX}{source}{}{sourcepointer-5}{}}
\newlabel{stacks-morphisms-lemma-composition-flat@cref}{{[originalreference][0][]06PX}{[1][0][]source}{}{}{}}
\makeatother
\crefname{originalreference}{original reference}{original references}
\Crefname{originalreference}{Original reference}{Original references}
\makeatletter
\newlabel{stacks-morphisms-lemma-flat-permanence}{{06Q0}{source}{}{sourcepointer-6}{}}
\newlabel{stacks-morphisms-lemma-flat-permanence@cref}{{[originalreference][0][]06Q0}{[1][0][]source}{}{}{}}
\makeatother
\crefname{originallemma}{original lemma}{original lemmas}
\Crefname{originallemma}{Original lemma}{Original lemmas}
\makeatletter
\newlabel{lemma-versal-ring-field-extension}{{0DQY}{source}{}{sourcepointer-7}{}}
\newlabel{lemma-versal-ring-field-extension@cref}{{[originallemma][0][]0DQY}{[1][0][]source}{}{}{}}
\makeatother
\crefname{originalreference}{original reference}{original references}
\Crefname{originalreference}{Original reference}{Original references}
\makeatletter
\newlabel{spaces-morphisms-section-points-finite-type}{{06EE}{source}{}{sourcepointer-8}{}}
\newlabel{spaces-morphisms-section-points-finite-type@cref}{{[originalreference][0][]06EE}{[1][0][]source}{}{}{}}
\makeatother
\crefname{originalreference}{original reference}{original references}
\Crefname{originalreference}{Original reference}{Original references}
\makeatletter
\newlabel{morphisms-section-points-finite-type}{{01T9}{source}{}{sourcepointer-9}{}}
\newlabel{morphisms-section-points-finite-type@cref}{{[originalreference][0][]01T9}{[1][0][]source}{}{}{}}
\makeatother
\crefname{originalreference}{original reference}{original references}
\Crefname{originalreference}{Original reference}{Original references}
\makeatletter
\newlabel{algebra-lemma-completion-flat}{{00MB}{source}{}{sourcepointer-10}{}}
\newlabel{algebra-lemma-completion-flat@cref}{{[originalreference][0][]00MB}{[1][0][]source}{}{}{}}
\makeatother
\title{The morphism from a versal ring to the stack is flat}
\author{The Stacks Project Authors}
\date{stacks-project:stacks-geometry.tex:local-snapshot-20261012; GNU FDL 1.2 or later}
\begin{document}
\maketitle
\noindent\textit{Editorial note.} All delivered mathematical text is the author's own, copied verbatim from the pinned Stacks Project chapter file stacks-geometry.tex; no statement or proof step was rewritten, reordered or removed.\par\medskip
\noindent\textit{Editorial note.} Span recomputation by label pairing (LANE\_BRIEF 10.2): in this chapter the author's proof environment follows the statement environment as a standalone block; the spec covers the complete environment pair (statement span, proof span) as recorded in the pool candidate.\par\medskip
\noindent\textit{Editorial note.} Occupation scan: combined claim\_key + statement\_anchor + generic-line scan over all registration files (active and withdrawn) found 0 hits for stacks-geometry.tex (ge-occupied.json, 2026-10-14); 207 route=geometry / chapter-list field hits in focused-seeds are alias noise (route name and alphabetical chapter enumerations), excluded per the ladder alias-blind-spot caveat.\par\medskip
\ifdefined\hypertarget\hypertarget{sourcepointer-1}{}\fi
\noindent\textit{Source reference.} Original reference 0DQU: sources/stacks-chow/tags.txt, source lines 17199--17199. Cross-chapter standard prerequisite cited by the author inside the counted span; it is used as a pointer and its own proof is not part of this item.\par\medskip
\ifdefined\hypertarget\hypertarget{sourcepointer-2}{}\fi
\noindent\textit{Source reference.} Original remark 0DR1: sources/stacks-chow/tags.txt, source lines 17206--17206. Cross-chapter standard prerequisite cited by the author inside the counted span; it is used as a pointer and its own proof is not part of this item.\par\medskip
\ifdefined\hypertarget\hypertarget{sourcepointer-3}{}\fi
\noindent\textit{Source reference.} Original lemma 0DR0: sources/stacks-chow/tags.txt, source lines 17205--17205. Cross-chapter standard prerequisite cited by the author inside the counted span; it is used as a pointer and its own proof is not part of this item.\par\medskip
\ifdefined\hypertarget\hypertarget{sourcepointer-4}{}\fi
\noindent\textit{Source reference.} Original lemma 0DQX: sources/stacks-chow/tags.txt, source lines 17202--17202. Cross-chapter standard prerequisite cited by the author inside the counted span; it is used as a pointer and its own proof is not part of this item.\par\medskip
\ifdefined\hypertarget\hypertarget{sourcepointer-5}{}\fi
\noindent\textit{Source reference.} Original reference 06PX: sources/stacks-chow/tags.txt, source lines 8586--8586. Cross-chapter standard prerequisite cited by the author inside the counted span; it is used as a pointer and its own proof is not part of this item.\par\medskip
\ifdefined\hypertarget\hypertarget{sourcepointer-6}{}\fi
\noindent\textit{Source reference.} Original reference 06Q0: sources/stacks-chow/tags.txt, source lines 8589--8589. Cross-chapter standard prerequisite cited by the author inside the counted span; it is used as a pointer and its own proof is not part of this item.\par\medskip
\ifdefined\hypertarget\hypertarget{sourcepointer-7}{}\fi
\noindent\textit{Source reference.} Original lemma 0DQY: sources/stacks-chow/tags.txt, source lines 17203--17203. Cross-chapter standard prerequisite cited by the author inside the counted span; it is used as a pointer and its own proof is not part of this item.\par\medskip
\ifdefined\hypertarget\hypertarget{sourcepointer-8}{}\fi
\noindent\textit{Source reference.} Original reference 06EE: sources/stacks-chow/tags.txt, source lines 8213--8213. Cross-chapter standard prerequisite cited by the author inside the counted span; it is used as a pointer and its own proof is not part of this item.\par\medskip
\ifdefined\hypertarget\hypertarget{sourcepointer-9}{}\fi
\noindent\textit{Source reference.} Original reference 01T9: sources/stacks-chow/tags.txt, source lines 2366--2366. Cross-chapter standard prerequisite cited by the author inside the counted span; it is used as a pointer and its own proof is not part of this item.\par\medskip
\ifdefined\hypertarget\hypertarget{sourcepointer-10}{}\fi
\noindent\textit{Source reference.} Original reference 00MB: sources/stacks-chow/tags.txt, source lines 817--817. Cross-chapter standard prerequisite cited by the author inside the counted span; it is used as a pointer and its own proof is not part of this item.\par\medskip
\begin{lemma}
\label{lemma-versal-ring-flat}
In Situation \ref{situation-versal} let $x_0 : \Spec(k) \to \mathcal{X}$
be a morphism, where $k$ is a finite type field over $S$.
Let $A$ be a versal ring to $\mathcal{X}$ at $x_0$.
Then the morphism $\Spec(A) \to \mathcal{X}$ of
Remark \ref{remark-upgrade} is flat.
\end{lemma}

\begin{proof}
If the local ring of $S$ at the image point is a G-ring, then this
follows immediately from
Lemma \ref{lemma-Artin-approximation-by-smooth-morphism}
and the fact that the map from a Noetherian local ring to its
completion is flat. In general we prove it as follows.

\medskip\noindent
Step I. If $A$ and $A'$ are two versal rings to $\mathcal{X}$ at $x_0$,
then the result is true for $A$ if and only if it is true for $A'$.
Namely, after possible swapping $A$ and $A'$, we may assume there is
a formally smooth map $\varphi : A \to A'$ such that the composition
$$
\Spec(A') \to \Spec(A) \to \mathcal{X}
$$
is the morphism $\Spec(A') \to \mathcal{X}$, see
Lemma \ref{lemma-versal-ring} and
Remark \ref{remark-upgrade}.
Since $A \to A'$ is faithfully flat we obtain the equivalence from
Morphisms of Stacks, Lemmas \ref{stacks-morphisms-lemma-composition-flat} and
\ref{stacks-morphisms-lemma-flat-permanence}.

\medskip\noindent
Step II. Let $l/k$ be a finite extension of fields. Let
$x_{l, 0} : \Spec(l) \to \mathcal{X}$ be the induced morphism.
Let $A$ be a versal ring to $\mathcal{X}$ at $x_0$ and
let $A \to A'$ be as in Lemma \ref{lemma-versal-ring-field-extension}.
Then again the composition
$$
\Spec(A') \to \Spec(A) \to \mathcal{X}
$$
is the morphism $\Spec(A') \to \mathcal{X}$, see Remark \ref{remark-upgrade}.
Arguing as before and using step I to see choice of versal rings
is irrelevant, we see that the lemma holds for $x_0$ if and only
if it holds for $x_{l, 0}$.

\medskip\noindent
Step III. Choose a scheme $U$ and a surjective smooth morphism
$U \to \mathcal{X}$. Then we can choose a finite type point
$z_0$ on $Z = U \times_\mathcal{X} x_0$ (this is a nonempty
algebraic space). Let $u_0 \in U$ be the image of $z_0$ in $U$.
Choose a scheme and a surjective \'etale map $W \to Z$
such that $z_0$ is the image of a closed point $w_0 \in W$
(see Morphisms of Spaces, Section
\ref{spaces-morphisms-section-points-finite-type}).
Since $W \to \Spec(k)$ and $W \to U$ are of finite type,
we see that $\kappa(w_0)/k$ and $\kappa(w_0)/\kappa(u_0)$
are finite extensions of fields
(see Morphisms, Section \ref{morphisms-section-points-finite-type}).
Applying Step II twice we may replace $x_0$ by
$u_0 \to U \to \mathcal{X}$. Then we see our morphism is the composition
$$
\Spec(\mathcal{O}_{U, u_0}^\wedge) \to U \to \mathcal{X}
$$
The first arrow is flat because completion of Noetherian local rings
are flat (Algebra, Lemma \ref{algebra-lemma-completion-flat})
and the second arrow is flat as
a smooth morphism is flat. The composition is flat as composition
preserves flatness.
\end{proof}

\end{document}
